% Outline for a numerical methods paper on the monolithic beam + rigid body
% solver. LaTeX version of paperOutline.txt.
%
% Each section lists what should go in it. The [material: ...] notes name
% the file in documentation/ that already holds the material (see the table
% in the "Material" section). For submission, swap the document class for
% the journal's (e.g. elsarticle for CMAME / Comput. Struct., or the Wiley
% NME template) and replace the itemised notes with text.

\documentclass[11pt,a4paper]{article}

\usepackage[T1]{fontenc}
\usepackage[utf8]{inputenc}
\usepackage[margin=25mm]{geometry}
\usepackage{amsmath,amssymb}
\usepackage{booktabs}
\usepackage{array}
\usepackage{enumitem}
\usepackage[hidelinks]{hyperref}

\setlist{itemsep=2pt,topsep=4pt}

% Vectors and tensors
\newcommand{\vect}[1]{\boldsymbol{#1}}
\newcommand{\tens}[1]{\boldsymbol{#1}}
\newcommand{\skewm}[1]{[\vect{#1}]_{\times}}
\newcommand{\dt}{\Delta t}

% Pointer to the documentation file that holds the material
\newcommand{\material}[1]{\par\smallskip{\raggedright\noindent\textit{[material: #1]}\par}}
\newcommand{\file}[1]{\texttt{#1}}

\title{A monolithic Newton--Raphson coupling of a finite-volume geometrically
  exact beam and a rigid body, for moored floating bodies\\[4pt]
  \large Paper outline}
\author{Colm Mc Alister}
\date{\today}

\begin{document}

\maketitle

\section*{About this outline}

This is a suggested layout for writing up the monolithic beam + rigid body
coupling as a paper for a numerical methods journal (for example
\emph{Int.\ J.\ Numer.\ Meth.\ Engng}, \emph{Comput.\ Methods Appl.\ Mech.\
Engrg}, \emph{J.\ Comput.\ Phys.}, \emph{Computers \& Structures}). For each
section it says what should go in it, and which file in
\file{documentation/} already holds the material.

Literature references are not listed here: each claim about prior work needs
a source checked by hand before it goes in.

\paragraph{Alternative titles.}
\begin{itemize}
  \item Monolithic finite-volume beam--rigid body coupling with Newmark
        integration in displacement form
  \item Removing the partitioned stability limit in beam-moored rigid body
        simulations
\end{itemize}

\paragraph{Material files} (all in \file{documentation/}):

\begin{center}
\begin{tabular}{@{}l>{\raggedright\arraybackslash}p{0.85\textwidth}@{}}
  \toprule
  Key & File \\
  \midrule
  CP & \file{monolithicRigidBodyCodePath.txt} \\
  PC & \file{pimpleToRigidBodySolverCallChain.txt} \\
  ML & \file{monolithicMatrixLayout.txt} \\
  TC & \file{twoCellMonolithicMatrix.txt} \\
  NB & \file{newmarkAndMonolithicBodyRow.txt} \\
  OC & \file{outerCorrectorsAndNewmark.pdf}; sources
       \file{source/outerCorrectorsAndNewmark.html}, \file{source/model.py},
       \file{source/order.py} \\
  MP & \file{monolithicCouplingPlan.pdf} \\
  TR & \file{monolithicSolverTestProgressReport.pdf} \\
  P5 & \file{phase5MultiLinePlan.md} \\
  \bottomrule
\end{tabular}
\end{center}

\section*{Abstract (150--250 words)}

\begin{itemize}
  \item Problem: a floating rigid body moored by flexible lines, inside a CFD
        simulation. The body and lines are usually coupled by partitioned
        iteration.
  \item Difficulty: partitioned coupling becomes unstable for light bodies and
        stiff, short lines (the limit
        $\beta\,\dt^2 EA/(L\,m_\mathrm{pin}) \gtrsim 1$).
  \item Method: the body's Newton--Euler equations, integrated with
        Newmark-$\beta$ in displacement form, are added as 6 extra rows and
        columns to the finite-volume beam's Newton system. The joint
        compatibility is eliminated, so beam and body are advanced together
        by one Newton iteration.
  \item Results: agrees with analytic and partitioned reference solutions;
        second order in time; stable where the partitioned scheme diverges;
        one PIMPLE outer corrector is enough, which reduces cost.
  \item Implementation in OpenFOAM (beamFoam / moorFV), with a single switch
        between partitioned and monolithic coupling.
\end{itemize}

\section{Introduction}

\begin{itemize}
  \item Motivation: floating offshore structures (wave energy converters,
        floating wind, aquaculture) moored by lines that bend, twist and go
        slack. CFD of the body needs a mooring model that is both accurate
        and robust.
  \item Mooring models: quasi-static, lumped-mass (e.g.\ MoorDyn), finite
        element and finite volume beam models.
  \item Coupling: partitioned (body solver and line solver take turns) versus
        monolithic (one system). Partitioned stability problems for light
        bodies and stiff connections (added-mass-like instability).
  \item Gap: no finite-volume geometrically exact beam with the rigid body in
        the same Newton system (check this claim against the literature).
  \item Contributions:
  \begin{enumerate}
    \item Monolithic formulation: body residuals in force units,
          Newmark-$\beta$ in displacement form, rotations by a step rotation
          vector.
    \item Elimination of the joint (pin and welded) so the attachment
          displacement is not an unknown.
    \item Block structure of the coupled Jacobian, and its physical
          interpretation (action--reaction in the columns, non-symmetry).
    \item Stability analysis of the partitioned scheme on a linear model,
          giving the limit the monolithic scheme removes.
    \item Verification and application to moored floating bodies in waves.
  \end{enumerate}
  \item Paper layout paragraph.
\end{itemize}

\section{Governing equations}

\subsection{Geometrically exact beam}
\begin{itemize}
  \item Total Lagrangian, Simo--Reissner type: displacement $\vect{W}$,
        rotation tensor $\tens{\Lambda}$ (incremental rotation
        $\vect{\Theta}$), internal force $\vect{Q}$ and moment $\vect{M}$.
  \item Balance of linear and angular momentum per unit length, including
        the lever-arm term $\vect{R}' \times \vect{Q}$.
  \item Constitutive law in section axes:
        $\tens{C}_Q = \mathrm{diag}(EA, GA, GA)$,
        $\tens{C}_M = \mathrm{diag}(GJ, EI_{yy}, EI_{zz})$.
        $J$ is the torsion constant ($\pi R^4/2$ for a solid circle).
  \item Loads: gravity, distributed hydrodynamic load (if any), ground
        contact.
\end{itemize}
\material{ML section 9 (coefficient origins), TC ``What each symbol is''}

\subsection{Rigid body}
\begin{itemize}
  \item Newton--Euler about the centre of mass:
  \begin{align*}
    m\,\vect{a} + c\,\vect{v} &= \vect{F}_\mathrm{ext} + \vect{F}_\mathrm{beam},\\
    \tens{J}\vect{\alpha} + \vect{\omega}\times\tens{J}\vect{\omega}
      + c_r\,\vect{\omega} &= \vect{M}_\mathrm{ext}
      + \vect{r}\times\vect{F}_\mathrm{beam}\;(+\,\vect{M}_\mathrm{beam}).
  \end{align*}
  \item Orientation $\tens{R} = \exp(\skewm{\theta})\,\tens{R}_0$, inertia
        $\tens{J} = \tens{R}\,\tens{J}_\mathrm{body}\,\tens{R}^T$.
  \item Fluid force and moment from the CFD as external loads; motion
        constraints (e.g.\ surge--heave--pitch only) as projections.
\end{itemize}
\material{ML 9.4, CP section 4}

\subsection{Joint between beam end and body}
\begin{itemize}
  \item Compatibility: attachment point
        $\vect{u}_b = \vect{x} + (\tens{R} - \tens{I})\,\vect{r}_0$.
  \item Pin: force passes, no moment (zero-moment end condition).
  \item Welded: end section rotates with the body; force and moment pass.
  \item Action--reaction: the body receives $-\vect{Q}_b$ (and $-\vect{M}_b$
        if welded).
\end{itemize}
\material{TC (layout, $\tens{B}_b$), ML section 6}

\section{Discretisation}

\subsection{Finite-volume beam in space}
\begin{itemize}
  \item Cell-centred unknowns, face fluxes of force and moment.
  \item Newton linearisation of the face force and moment: the face
        coefficients $\tens{k}$, $\tens{\kappa}$, $\tens{g}$,
        $\tens{\gamma}$, $\tens{K}^M$, $\tens{\eta}$ (definitions and
        physical meaning). Boundary faces, half-cell distance
        $\Delta_p$ (\file{pDelta}).
  \item Pin end factor
        $\tens{B} = (\tens{\kappa} + \tens{\gamma})^{-1}\tens{\kappa}$:
        the end face rotation eliminated by the zero-moment condition.
\end{itemize}
\material{ML sections 9.1--9.3, TC ``What each symbol is'' and ``Notation''}

\subsection{Time integration}
\begin{itemize}
  \item Beam: Newmark (as implemented) with its predictor.
  \item Body: Newmark-$\beta$ in displacement form, with predictor and
        acceleration
        \begin{equation*}
          \tilde{\vect{x}} = \vect{x}_n + \dt\,\vect{v}_n
            + \dt^2\left(\tfrac12 - \beta\right)\vect{a}_n,
          \qquad
          \vect{a}_{n+1} = \frac{\vect{x}_{n+1} - \tilde{\vect{x}}}{\beta\,\dt^2},
        \end{equation*}
        and the same for the step rotation vector $\vect{\theta}$. Explain
        why the predictor has no $\beta\,\dt^2\,\vect{a}_{n+1}$ term.
  \item Rotation Jacobian $\tens{T}(\vect{\theta})$ and its closed form.
\end{itemize}
\material{ML sections 9.4 and 10, NB sections 1 and 1a}

\section{Monolithic coupling}

The core of the paper.

\subsection{Unknowns and residuals}
\begin{itemize}
  \item Unknowns $[\Delta\vect{W}, \Delta\vect{\Theta}]$ per cell and
        $[\Delta\vect{x}, \Delta\vect{\theta}]$ for the body, in
        $6\times 6$ blocks.
  \item Body residuals in force units:
  \begin{align*}
    \vect{R}_x &= \frac{m\,(\vect{x} - \tilde{\vect{x}})}{\beta\,\dt^2}
      + c\,\vect{v} - \vect{F}_\mathrm{ext} + \vect{Q}_b,\\
    \vect{R}_\theta &= \tens{J}\vect{\alpha}
      + \vect{\omega}\times\tens{J}\vect{\omega} + c_r\,\vect{\omega}
      - \vect{M}_\mathrm{ext} + \vect{r}\times\vect{Q}_b.
  \end{align*}
  \item Statement and short proof that a converged Newton iterate satisfies
        the Newmark-$\beta$ update exactly (the linearisation of
        $\vect{Q}_b$ is the only approximation per iteration).
\end{itemize}
\material{NB sections 2--4}

\subsection{Eliminating the joint}
\begin{itemize}
  \item $\Delta\vect{u}_b = \Delta\vect{x} - \tens{A}\,\Delta\vect{\theta}$,
        with $\tens{A} = \skewm{r}\,\tens{T}$.
  \item How the attachment-face source term
        $\tens{K}_b(\vect{W}_b - \vect{W}_{b,\mathrm{prev}})$ becomes the
        body columns of the matrix.
\end{itemize}
\material{MP section 2, ML section 6}

\subsection{Linearisation and block structure}
\begin{itemize}
  \item The coupling blocks (table), and the body diagonal including the
        turning-arm term $\skewm{Q_b}\tens{A}$ and the gyroscopic terms.
  \item Physical reading: every coupling entry is the linearised face force
        $\Delta\vect{Q}_b$ passed through a row operator ($\tens{I}$,
        $\tens{I}$, lever arm $\tens{K}^M$, $\skewm{r}$). Action--reaction
        appears down the columns, not across transposes.
  \item Why the Jacobian is not symmetric even with consistent signs.
  \item Pin versus welded: extra blocks for the welded joint.
  \item Motion constraints as projections of the body rows and columns.
\end{itemize}
\material{ML sections 1--8, TC ``Where the action--reaction balance is''}

\subsection{Solution of the linear system}
\begin{itemize}
  \item Assembly into one sparse matrix; sparse LU (non-symmetric).
  \item Sparsity: block-tridiagonal beam plus one bordering block row and
        column at the attachment cell. Cost relative to the beam alone.
\end{itemize}
\material{ML sections 3--4, CP sections 5--6}

\subsection{Algorithm}
\begin{itemize}
  \item Pseudo-code for one time step: new time step (store state, reset
        $\vect{\theta}$), Newton loop (assemble beam, add body blocks, solve,
        update beam and body, check convergence), accept (update
        $\vect{v}$, $\vect{a}$, $\vect{\omega}$, $\vect{\alpha}$).
\end{itemize}
\material{CP, NB section 3}

\section{The partitioned scheme and its stability limit}

\begin{itemize}
  \item The partitioned algorithm used for comparison
        (\file{FvBeamNewmark}): body Newmark step against the frozen beam
        force, then the beam with the end prescribed, repeated in the outer
        correctors.
  \item Linear body--spring model: amplification matrices for the
        partitioned loop and the monolithic update. Show the partitioned
        iteration diverges when $\beta\,\dt^2 k/m_\mathrm{pin} \gtrsim 1$,
        with $m_\mathrm{pin} = 1/(1/m + a^2/J)$ for an off-centre
        attachment.
  \item One outer corrector: partitioned is explicit, first order and
        unstable; monolithic is implicit and second order.
\end{itemize}
\material{OC (figures \file{growth.png}, \file{newton.png}),
  \file{source/model.py}, \file{source/order.py}, the
  \file{rigidBodyEndWhirl} sweeps}

\section{Coupling with the flow solver}

Keep this short: it is context, not the contribution.
\begin{itemize}
  \item interFoam PIMPLE loop, mesh motion with
        \file{moveMeshOuterCorrectors}.
  \item Fluid force held fixed during the beam + body Newton loop and
        updated at each outer corrector (the fluid stays partitioned).
  \item Run-time selection of partitioned or monolithic coupling.
\end{itemize}
\material{PC, OC sections 1--2}

\section{Verification}

Each test: setup, what it checks, the result, a figure or table.
\begin{itemize}
  \item Translation only: body on a beam spring, against an analytic
        spring--mass solution.
  \item Rotation: off-centre attachment and twist
        (\file{rigidBodyEndTwist}), pin and welded joints.
  \item Order of accuracy in time (expect second order), for both
        translation and rotation.
  \item Newton convergence: quadratic convergence of the residual norm
        (show the exact Jacobian pays off).
  \item Monolithic vs converged partitioned on cases where both run: they
        should agree to the iteration tolerance.
  \item Robustness: \file{rigidBodyEndWhirl} sweeps of moment and $EI$; the
        partitioned limit seen numerically, matching the stability section.
\end{itemize}
\material{TR, \file{tutorials/rigidBodyEnd*}, \file{rigidBodyAndBeam\_loop}
  and \file{\_beamFoamCoupled}, the regression check of 2026-10-20
  (agreement to $\sim 2\times10^{-6}$ with identical Newton counts after the
  legacy removal)}

\section{Application: moored floating bodies in waves}

\begin{itemize}
  \item \file{floatRigidBody}: floating box with a mooring line in waves;
        motion and line force, monolithic vs partitioned vs MoorDyn.
  \item Slack line (\file{floatRigidBody\_slackLine}): snap loads and
        compression, where line bending and ground contact matter.
  \item Two lines (\file{floatRigidBody\_twoLine}), if Phase 5 is complete.
  \item Sweeps of line stiffness and wave period
        (\file{floatRigidBody\_sweeps}): where the partitioned scheme fails,
        and cost.
  \item Cost: Newton iterations per step (\file{beamSolverIterations}
        function object), outer correctors needed, wall time split fluid /
        beam.
\end{itemize}
\material{\file{tutorials/floatRigidBody*},
  \file{plot\_floatRigidBody\_sweeps.py},
  \file{plot\_monolithic\_vs\_partitioned\_*.py}}

\section{Discussion}

\begin{itemize}
  \item When monolithic coupling matters (the stability number) and when
        partitioned is good enough.
  \item Cost: the extra rows are cheap; the saving comes from fewer outer
        correctors and fewer failed runs.
  \item Limitations:
  \begin{itemize}
    \item one line per body today (Phase 5 plan for several lines);
    \item the fluid stays partitioned (no fluid added mass in the monolithic
          system);
    \item non-symmetric Jacobian, so a direct sparse LU; scaling to long
          lines and many lines;
    \item geometric terms that are only first order in $\vect{\theta}$
          ($\tens{T} = \tens{I}$ limit).
  \end{itemize}
  \item Extensions: several lines on one body, several bodies, implicit
        fluid added mass, iterative solvers with block preconditioners.
\end{itemize}
\material{P5, OC}

\section{Conclusions}

\begin{itemize}
  \item Restate the contributions and the main quantitative results (the
        stability number, order of accuracy, cost).
\end{itemize}

\appendix

\section{The full two-cell matrix in physical terms}
With the meaning of every entry.
\material{TC}

\section{Derivation of the face coefficients}
$\tens{k}$, $\tens{\kappa}$, $\tens{g}$, $\tens{\gamma}$, $\tens{K}^M$,
$\tens{\eta}$, and the pin factor $\tens{B}$.
\material{ML section 9, TC ``Notation''}

\section{Rotation}
Exponential map, rotation Jacobian $\tens{T}(\vect{\theta})$, and the
rotational inertia Jacobian.
\material{ML 9.4}

\section{Linear stability model}
Of the partitioned and monolithic schemes.
\material{OC, \file{source/model.py}}

\section{Implementation notes and code availability}
Repositories, branch, and the tutorial cases that reproduce each figure.
\material{CP, PC}

\section*{Figures and tables to prepare}

\begin{center}
\begin{tabular}{@{}lp{0.8\textwidth}@{}}
  \toprule
  Item & Content \\
  \midrule
  Fig.~1 & Sketch: body, joint, beam cells and faces, with $\vect{r}$,
           $\vect{Q}_b$, $\vect{x}$, $\vect{\theta}$. \\
  Fig.~2 & Sparsity pattern of the coupled Jacobian (beam blocks + bordered
           body block). \\
  Fig.~3 & Linear model: spectral radius vs
           $\beta\,\dt^2 k/m_\mathrm{pin}$, partitioned and monolithic
           (\file{growth.png}). \\
  Fig.~4 & Order of accuracy in time. \\
  Fig.~5 & Newton residual vs iteration (quadratic convergence). \\
  Fig.~6 & Verification time histories against analytic / reference. \\
  Fig.~7 & Robustness map: where partitioned diverges, monolithic runs. \\
  Fig.~8 & \file{floatRigidBody} motion and line force: monolithic,
           partitioned, MoorDyn. \\
  Fig.~9 & Cost: Newton iterations and wall time per step, sweeps. \\
  Tab.~1 & Notation. \\
  Tab.~2 & Coupling blocks and their physical meaning. \\
  Tab.~3 & Case parameters. \\
  \bottomrule
\end{tabular}
\end{center}

\section*{Gaps to fill before writing}

\begin{itemize}
  \item A clean order-of-accuracy study in the full code (not only the
        linear model), for translation and rotation.
  \item A Newton convergence plot from the code showing quadratic
        convergence, with and without the welded terms.
  \item Wall-time measurements on the same machine for monolithic vs
        partitioned with the outer corrector count each needs.
  \item A literature review to support the ``gap'' claim in the
        introduction and to place the method against monolithic FSI and
        mooring coupling work.
\end{itemize}

\end{document}
